\documentclass{article}
\usepackage{../header}
\title{Automata and Formal Languages}
\author{Notes made by Finley Cooper}

\begin{document}
  \maketitle
  \newpage
  \tableofcontents
  \newpage
  \section{Introduction}
  In this course we will have $ 0\in \N $.
\begin{definition}
	(Alphabet) An \textit{alphabet} $ \Sigma $ is a finite non-empty set with elements which are \textit{symbols} or \textit{letters}.
\end{definition}
\begin{definition}
	(String) A \textit{string} or \textit{word} of length $ m $ over $ \Sigma $ is $ \alpha = x_1x_2\dots x_m $ with $ x_i\in \Sigma $ or more formally $ \alpha \in\Sigma^m $. We write $ |\alpha| = m $ for the length of $ \alpha $.
\end{definition}
\begin{remark}
  We will denote the empty word as $ \varepsilon $ and denote
  \[
	  \Sigma^* = \bigcup_{n\in \N}\Sigma^n
  \]
  and $ \Sigma^\dagger = \Sigma^* \setminus \{\varepsilon\} $.
\end{remark}
\begin{definition}
	(Concatenation) If $ \alpha = x_1\cdots x_m\in \Sigma^m $ and $ \beta= y_1\cdots y_n\in \Sigma^n $ then $ \alpha\beta $ is $ x_1\cdots x_my_1\cdots y_n \in \Sigma^{m+n} $.
\end{definition}
For $ 0\le i \le j \le m $, $ x_{i+1}\cdots x_j $ is a substring of $ \alpha $ of length $ j-i $. So for $ \alpha,\beta\in \Sigma^* $, $\beta $ is a substring of $ \alpha $ if and only if there exists $ \gamma,\delta\in\Sigma^* $ with $ \alpha = \gamma\beta\delta $.
\begin{definition}
	(Homomorphism) For alphabets $ \Sigma, \Pi $, a \textit{homomorphism} $ h: \Sigma^* \to \Pi^* $ satisfies $ h(\alpha\beta) = h(\alpha)h(\beta) $ for all $ \alpha,\beta\in \Sigma^* $.
\end{definition}
\begin{remark}
  This implies that $ h(\varepsilon) = \varepsilon $. 
\end{remark}
\begin{remark}
  Alternatively we can just consider a function $ f : \Sigma \to \Pi$ which defines a unique homomorphism extending $ f $ by concatenation.
\end{remark}
\begin{definition}
	(Formal Langauge) A \textit{formal language} $ L $ over $ \Sigma $ is \textit{any} subset of $ \Sigma^* $.
\end{definition}
\section{Regular Languages and Finite-State Automata}
\subsection{Deterministic Finite-State Automata}
\begin{remark}
	Suppose we want to describe a language $ L $ over $ \Sigma $ using a finite amount of information. We could extend $ \Sigma $ with other symbols to get an alphabet $ \Pi \supset \Sigma $. and assign a word $ w\in\Pi^* $ to $ L $.
\end{remark}
We can do this for countably many $ L $. But we want $ w $ to \textit{describe} $ L $. We'll now see a countable class of languages where we can do this symbolically. We first give a computational method to determine these languages.
\begin{definition}
	(Deterministic Finite-State Automata) A \textit{Deterministic Finite-State Automata} (DFA) is a structure
	\[
	  D = (Q,\Sigma, \delta, q_0, F)
	\]
	where,
	\begin{enumerate}
		\item $ Q $ is a finite non-empty set of states;
		\item $ \Sigma $ is an alphabet;
		\item $ \delta $ is a transition function from $ Q\times\Sigma \to Q $;
		\item Starting state $ q_0\in Q $;
		\item $ F\subseteq Q $ is a set of acceptance states.
	\end{enumerate}
\end{definition}
The idea for a DFA is that we input a word $ w\in \Sigma^* $ we call $ w = x_1\cdots x_n $. The DFA starts in state $ q_0 $ and reads the first letter of the input string $ x_1 $. The DFA then transitions according to $ \delta $ to $ q_1 = \delta(q_0,x_1)\in Q $. We then repeat this process with $ x_2 $ and so on. Once the string is exhusted the string is accepted if $ q_n\in F $ otherwise it is rejected. This defines a language.

  \end{document}
